% Generated draft: requires mathematical and source review.
\documentclass[11pt]{article}
\usepackage[margin=25mm]{geometry}
\usepackage{fontspec}
\setmainfont{Latin Modern Roman}
\usepackage{amsmath,amssymb,mathtools,unicode-math}
\setmathfont{Latin Modern Math}
\usepackage{enumitem,needspace,xurl,hyperref}
\setlength{\parindent}{0pt}
\setlength{\parskip}{5pt}
\setlength{\emergencystretch}{3em}
\newcommand{\N}{\mathbb{N}}
\newcommand{\Z}{\mathbb{Z}}
\newcommand{\Q}{\mathbb{Q}}
\newcommand{\R}{\mathbb{R}}
\newcommand{\C}{\mathbb{C}}
\newcommand{\F}{\mathbb{F}}
\newcommand{\A}{\mathbb{A}}
\newcommand{\PP}{\mathbb{P}}
\newcommand{\E}{\mathbb{E}}
\newcommand{\eps}{\varepsilon}
\newcommand{\dd}{\,\mathrm{d}}
\begin{document}
\section*{UnsolvedMath ID: 11000312}
\textbf{Source problem code:} catalogue entry
\subsection{Definitions and mathematical statement}
Let $\operatorname{IOut}_n=\ker(\operatorname{Out}(F_n)\to\operatorname{GL}_n(\Z))$ in the customary convention. The quotient acts on its rational cohomology by conjugation, well defined because inner conjugations act trivially on cohomology; write $(-)^{\operatorname{GL}}$ for the invariant part under this action.

Problem 4.5 asks to prove or disprove that restriction is an isomorphism
\[
 r:H^4(\operatorname{Out}(F_4);\Q)\cong\Q
 \longrightarrow
 H^4(\operatorname{IOut}_4;\Q)^{\operatorname{GL}},
\]
with the right-hand side generated by a nonzero multiple of the source class $\tau_2$. Equivalently, in this interpretation it asks that the target be one-dimensional and that a generator $u$ of the source restrict to $r(u)=c\tau_2$ for some $c\in\Q\setminus\{0\}$. The definition and normalization of $\tau_2$ are absent. The extracted degree placement in ``H 4'' is ambiguous, but the natural arrow to subgroup invariants is the cohomological restriction convention; ordinary homology has the opposite functorial direction.

The following passage sets up another specific problem without supplying its final question. It reports
\[
 H^8(M_{3,*};\Q)\cong\Q^2,\qquad
 H^8(\operatorname{GL}_6(\Z);\Q)\cong\Q,\qquad
 H^8(\operatorname{Out}(F_6);\Q)\cong\Q,
\]
and the natural maps (9)
\[
 M_{3,*}\xrightarrow{\ i\ }\operatorname{Out}(F_6)
 \xrightarrow{\ p\ }\operatorname{GL}_6(\Z),
\]
where $i$ is described as injective and $p$ is the homological projection. The meaning of $M_{3,*}$, the construction of $i$ and the degree-index conventions are not restated; a once-marked genus-three mapping class group is the usual possible interpretation, not a supplied definition. No rank, nonvanishing or isomorphism assertion about the induced degree-eight maps is invented.
\subsection{Short English statement}
Is degree-four restriction from $\operatorname{Out}(F_4)$ onto the general-linear invariants of its Torelli subgroup an isomorphism, generated by the named $\tau_2$ class? The class and notation conventions need confirmation.
\subsection{Sources}
\begin{itemize}
\item Local registry: \url{data/unsolvedmath/problems.json}, permanent UnsolvedMath ID $11000312$, exact source code \texttt{catalogue entry}.
\item Farb, \emph{Problems on Mapping Class Groups and Related Topics}, Problem 4.5, PDF page 366: \url{https://www.math.uchicago.edu/~farb/papers/mcgbook.pdf}. Supplied access date: 2026-07-29; extraction: pdf-font-structure; page-disambiguated numbering among 329 headings.
\item Source attributions for the degree-eight groups: Looijenga [61]; Elbaz-Vincent, Gangl and Soul\'e [13]; Ohashi [91]. Full entries, the definition of $\tau_2$ and the following specific question are absent. The registry cannot verify the modern status.
\end{itemize}
\end{document}
